    \section{Mandat no. 3 : Simulations
Monte-Carlo}\label{mandat-no.-3-simulations-monte-carlo}

    \subsection{Fonctions des temps
d'arrivés}\label{fonctions-des-temps-darrivuxe9s}

    \subsubsection{Processus de poisson et
exponentielle.}\label{processus-de-poisson-et-exponentielle.}

    La fonction de masse de la loi de Poisson est donnée par : \[
P(X=x)=\frac{e^{-\lambda T}(\lambda T)^x}{x!}
\] - Où \(\lambda\) est le taux d'arrivée. - T est une durée de temps.

Cette fonction donne la probabilité que \(x\) arrivées surviennent dans
un intervalle de temps \(T\). Cependant, ici on s'intéresse plutôt au
\emph{temps entre chaque arrivée}. Si \(x\) représente le nombre
d'arrivées dans l'intervalle \([0,T]\), posé \(x=0\) devrait donnée la
probabilité qu'aucune arrivée ne survienne dans cet intervalle.

Autrement dit, l'absence d'événement dans \([0,T]\) c'est la probabilité
que le temps jusqu'à la prochaine arrivée soit plus grand que \(T\): \[
P(T_A > T)
\]

\(x\) est donc remplacé par \(0\) dans la formule de Poisson, et c'est
simplifié: \[
P(X=0)=\frac{e^{-\lambda T}(\lambda T)^0}{0!}
\] \[
P(X=0)=\frac{e^{-\lambda T}\cdot 1}{1}
\] \[
P(X=0)=e^{-\lambda T}
\]

Ça donne une fonction exponentielle décroissante. Plus \(T\) augmente,
plus la probabilité diminue. C'est la fonction de survie: \[
P(T_A > T) = e^{-\lambda T}
\]

La CDF s'obtient en prenant le complément, étant donnée la présence d'un
\(>\) dans la formule de survie. \[
F(T)=P(T_A \le T)=1 - e^{-\lambda T}
\]

    \subsubsection{Inversion de la fonction}\label{inversion-de-la-fonction}

    Chaque joueur arrive à un temps \(P\), puis quitte après une session de
jeu de durée aléatoire \(Q\). Donc, chaque joueur est actif dans un
intervalle \([P, P+Q]\).

L'activité d'un joueur est donc: \[
P \le t \le P + Q
\]

Ça permet de déterminer combien de joueurs sont actifs à n'importe quel
instant (minute) \(t\). La génération des temps d'arrivée est fait avec
la méthode de transformation inverse. Il faut donc appliqué une
transformation inverse sur la CDF trouvé plus haut. On résoud pour
\(T\). \[
F(T)=1 - e^{-\lambda T}
\] \[
U = 1 - e^{-\lambda T}
\] \[
U - a = e^{-\lambda T}
\] \[
\ln(1 - U) = -\lambda T
\] \[
T = \frac{\ln(1 - U)}{-\lambda}
\] Ici, \(U\) est une variable aléatoire continue uniforme entre
\([0,1]\).

    \subsubsection{Python}\label{python}

La fonction python suivante génère les temps d'arrivés \(P\) en
utilisant la transformation inverse trouvé plus haut. Elle prend comme
paramètres le taux d'arrivés par minutes et le nombre de données à
générés.

    \begin{tcolorbox}[breakable, size=fbox, boxrule=1pt, pad at break*=1mm,colback=cellbackground, colframe=cellborder]
\begin{Verbatim}[commandchars=\\\{\}]
\PY{k}{def}\PY{+w}{ }\PY{n+nf}{generate\PYZus{}arrival\PYZus{}times}\PY{p}{(}\PY{n}{arrivals\PYZus{}per\PYZus{}minutes}\PY{p}{,} \PY{n}{sample\PYZus{}size}\PY{p}{)}\PY{p}{:}
    \PY{n}{uniform\PYZus{}randoms} \PY{o}{=} \PY{n}{np}\PY{o}{.}\PY{n}{random}\PY{o}{.}\PY{n}{rand}\PY{p}{(}\PY{n}{sample\PYZus{}size}\PY{p}{)}
    \PY{k}{return} \PY{o}{\PYZhy{}}\PY{n}{np}\PY{o}{.}\PY{n}{log}\PY{p}{(}\PY{n}{uniform\PYZus{}randoms}\PY{p}{)} \PY{o}{/} \PY{n}{arrivals\PYZus{}per\PYZus{}minutes}
\end{Verbatim}
\end{tcolorbox}

    \subsection{Temps de jeux des joueurs}\label{temps-de-jeux-des-joueurs}

La variable aléatoire normale \(Q\) a été déterminée dans le mandat 2,
avec sa moyenne et sa variance. Elle dicte le nombre d'heures qu'un
joueur passe sur le jeux.

La variable est indépendante du temps d'arrivé ou tout autres valeurs
associés au mandat-03.

Pour la générer l'implémentation manuel de l'algorithme de Box-Muller
est fesable, mais c'est inutile, puisque NumPy a déjà une implémentation
dans la fonction \texttt{randn}, qui à été confirmé comme équivalente
lors des laboratoires.

    \begin{tcolorbox}[breakable, size=fbox, boxrule=1pt, pad at break*=1mm,colback=cellbackground, colframe=cellborder]
\begin{Verbatim}[commandchars=\\\{\}]
\PY{k}{def}\PY{+w}{ }\PY{n+nf}{generate\PYZus{}play\PYZus{}durations}\PY{p}{(}\PY{n}{average}\PY{p}{,} \PY{n}{variance}\PY{p}{,} \PY{n}{sample\PYZus{}size}\PY{p}{)}\PY{p}{:}
    \PY{k}{return} \PY{n}{np}\PY{o}{.}\PY{n}{random}\PY{o}{.}\PY{n}{normal}\PY{p}{(}\PY{n}{average}\PY{p}{,} \PY{n}{variance}\PY{p}{,} \PY{n}{sample\PYZus{}size}\PY{p}{)}
\end{Verbatim}
\end{tcolorbox}

    \subsection{Paramètres de simulation}\label{paramuxe8tres-de-simulation}

    \begin{tcolorbox}[breakable, size=fbox, boxrule=1pt, pad at break*=1mm,colback=cellbackground, colframe=cellborder]
\begin{Verbatim}[commandchars=\\\{\}]
\PY{n}{tested\PYZus{}arrivals\PYZus{}per\PYZus{}minutes} \PY{o}{=} \PY{p}{[}\PY{l+m+mi}{10}\PY{p}{,} \PY{l+m+mi}{67}\PY{p}{,} \PY{l+m+mi}{100}\PY{p}{]} \PY{c+c1}{\PYZsh{} Taux d\PYZsq{}arrivés testés (branchements / minutes)}
\PY{n}{MU\PYZus{}Q} \PY{o}{=} \PY{l+m+mf}{280.56}                               \PY{c+c1}{\PYZsh{} Moyenne de temps de jeux d\PYZsq{}un joueur (minutes)}
\PY{n}{SIGMA\PYZus{}Q} \PY{o}{=} \PY{l+m+mf}{50.38}                             \PY{c+c1}{\PYZsh{} Variance du temps de jeux (minutes)}
\PY{n}{sample\PYZus{}size} \PY{o}{=} \PY{l+m+mi}{10000}                         \PY{c+c1}{\PYZsh{} N (minutes)}
\end{Verbatim}
\end{tcolorbox}

    \subsection{Fonction de simulation}\label{fonction-de-simulation}

Avec le taux de branchements, la transformation inverse calculé plus
haut est utilisé pour donnée un nombre de nouveau joueurs à chaque
minutes de la simulation. Leurs temps d'arrivé est calculé
progressivement au fur et à mesure que le temps avance. Par la suite, un
nombre de temps de jeux est généré pour chaque branchements. Le temps de
départ du joueur est ensuite calculé.

Le nombre de joueurs actif à n'importe quel instant \(t\) est finalement
calculé en appliquant \[
    T_{\text{arrivé}} \le t \le T_{\text{départ}}
\]

    \begin{tcolorbox}[breakable, size=fbox, boxrule=1pt, pad at break*=1mm,colback=cellbackground, colframe=cellborder]
\begin{Verbatim}[commandchars=\\\{\}]
\PY{k}{def}\PY{+w}{ }\PY{n+nf}{simulate}\PY{p}{(}\PY{n}{arrivals\PYZus{}per\PYZus{}minutes}\PY{p}{,} \PY{n}{simulation\PYZus{}time}\PY{p}{)}\PY{p}{:}
    \PY{n}{arrival\PYZus{}times} \PY{o}{=} \PY{p}{[}\PY{p}{]}
    \PY{n}{t} \PY{o}{=} \PY{l+m+mi}{0}

    \PY{c+c1}{\PYZsh{} Generating a P for each minute of the simulation}
    \PY{c+c1}{\PYZsh{} While making the arrivals later and later.}
    \PY{k}{while} \PY{n}{t} \PY{o}{\PYZlt{}} \PY{n}{simulation\PYZus{}time}\PY{p}{:}
        \PY{n}{t} \PY{o}{+}\PY{o}{=} \PY{n}{generate\PYZus{}arrival\PYZus{}times}\PY{p}{(}\PY{n}{arrivals\PYZus{}per\PYZus{}minutes}\PY{p}{,} \PY{l+m+mi}{1}\PY{p}{)}\PY{p}{[}\PY{l+m+mi}{0}\PY{p}{]}
        \PY{n}{arrival\PYZus{}times}\PY{o}{.}\PY{n}{append}\PY{p}{(}\PY{n}{t}\PY{p}{)}
    \PY{n}{arrival\PYZus{}times} \PY{o}{=} \PY{n}{np}\PY{o}{.}\PY{n}{array}\PY{p}{(}\PY{n}{arrival\PYZus{}times}\PY{p}{)}


    \PY{n}{play\PYZus{}time} \PY{o}{=} \PY{n}{generate\PYZus{}play\PYZus{}durations}\PY{p}{(}
        \PY{n}{MU\PYZus{}Q}\PY{p}{,}
        \PY{n}{SIGMA\PYZus{}Q}\PY{p}{,}
        \PY{n+nb}{len}\PY{p}{(}\PY{n}{arrival\PYZus{}times}\PY{p}{)}\PY{p}{)}
    \PY{n}{departures} \PY{o}{=} \PY{n}{arrival\PYZus{}times} \PY{o}{+} \PY{n}{play\PYZus{}time}

    \PY{c+c1}{\PYZsh{} This removes the warmup time of the population of servers.}
    \PY{c+c1}{\PYZsh{} Without this, the averages are off because they assume infinite}
    \PY{c+c1}{\PYZsh{} amount of time.}
    \PY{n}{t\PYZus{}start} \PY{o}{=} \PY{l+m+mf}{0.3} \PY{o}{*} \PY{n}{simulation\PYZus{}time}  \PY{c+c1}{\PYZsh{} remove transient phase}
    \PY{n}{t\PYZus{}values} \PY{o}{=} \PY{n}{np}\PY{o}{.}\PY{n}{linspace}\PY{p}{(}\PY{n}{t\PYZus{}start}\PY{p}{,} \PY{n}{simulation\PYZus{}time}\PY{p}{,} \PY{l+m+mi}{500}\PY{p}{)}

    \PY{n}{active\PYZus{}players} \PY{o}{=} \PY{p}{[}\PY{p}{]}
    \PY{k}{for} \PY{n}{t} \PY{o+ow}{in} \PY{n}{t\PYZus{}values}\PY{p}{:}
        \PY{n}{active} \PY{o}{=} \PY{n}{np}\PY{o}{.}\PY{n}{sum}\PY{p}{(}\PY{p}{(}\PY{n}{arrival\PYZus{}times} \PY{o}{\PYZlt{}}\PY{o}{=} \PY{n}{t}\PY{p}{)} \PY{o}{\PYZam{}} \PY{p}{(}\PY{n}{departures} \PY{o}{\PYZgt{}}\PY{o}{=} \PY{n}{t}\PY{p}{)}\PY{p}{)}
        \PY{n}{active\PYZus{}players}\PY{o}{.}\PY{n}{append}\PY{p}{(}\PY{n}{active}\PY{p}{)}

    \PY{k}{return} \PY{n}{arrival\PYZus{}times}\PY{p}{,} \PY{n}{play\PYZus{}time}\PY{p}{,} \PY{n}{t\PYZus{}values}\PY{p}{,} \PY{n}{active\PYZus{}players}
\end{Verbatim}
\end{tcolorbox}

    \subsection{Roulement de la
simulation}\label{roulement-de-la-simulation}

Le code suivant éxecute la simulation monte-carlo pour chaque taux de
branchements voulu. Elle vient sauvegarder certaines valeurs analysé
afin d'être utilisé dans les histogrammes et des affichages.

Elle prouve aussi la réponse au besoin d'être capable de dire combiens
de joueurs sont connecté à un instant \(t\), et fait une comparaison de
la moyenne théorique (\(\lambda\times\mu_Q\)) versus réel.

    \begin{tcolorbox}[breakable, size=fbox, boxrule=1pt, pad at break*=1mm,colback=cellbackground, colframe=cellborder]
\begin{Verbatim}[commandchars=\\\{\}]
\PY{n}{results} \PY{o}{=} \PY{p}{[}\PY{p}{]}

\PY{c+c1}{\PYZsh{} Testing every rates set in the test array up in this document}
\PY{k}{for} \PY{n}{rate} \PY{o+ow}{in} \PY{n}{tested\PYZus{}arrivals\PYZus{}per\PYZus{}minutes}\PY{p}{:}
    \PY{n}{arrivals}\PY{p}{,} \PY{n}{play\PYZus{}time}\PY{p}{,} \PY{n}{t}\PY{p}{,} \PY{n}{active\PYZus{}players} \PY{o}{=} \PY{n}{simulate}\PY{p}{(}\PY{n}{rate}\PY{p}{,} \PY{n}{sample\PYZus{}size}\PY{p}{)}
    \PY{n}{theorical\PYZus{}average} \PY{o}{=} \PY{n}{rate} \PY{o}{*} \PY{n}{MU\PYZus{}Q}
    \PY{n}{real\PYZus{}average} \PY{o}{=} \PY{n}{np}\PY{o}{.}\PY{n}{mean}\PY{p}{(}\PY{n}{active\PYZus{}players}\PY{p}{)}

    \PY{c+c1}{\PYZsh{} Points à différents instants}
    \PY{n}{t\PYZus{}check} \PY{o}{=} \PY{p}{[}\PY{l+m+mi}{10}\PY{p}{,} \PY{l+m+mi}{50}\PY{p}{,} \PY{l+m+mi}{500}\PY{p}{]}
    \PY{n}{active\PYZus{}at\PYZus{}t} \PY{o}{=} \PY{p}{\PYZob{}}\PY{n}{ti}\PY{p}{:} \PY{n}{np}\PY{o}{.}\PY{n}{sum}\PY{p}{(}\PY{p}{(}\PY{n}{arrivals} \PY{o}{\PYZlt{}}\PY{o}{=} \PY{n}{ti}\PY{p}{)} \PY{o}{\PYZam{}} \PY{p}{(}\PY{n}{arrivals} \PY{o}{+} \PY{n}{play\PYZus{}time} \PY{o}{\PYZgt{}}\PY{o}{=} \PY{n}{ti}\PY{p}{)}\PY{p}{)} \PY{k}{for} \PY{n}{ti} \PY{o+ow}{in} \PY{n}{t\PYZus{}check}\PY{p}{\PYZcb{}}

    \PY{c+c1}{\PYZsh{} Sauvegarde des résultats}
    \PY{n}{results}\PY{o}{.}\PY{n}{append}\PY{p}{(}\PY{p}{\PYZob{}}
        \PY{l+s+s2}{\PYZdq{}}\PY{l+s+s2}{rate}\PY{l+s+s2}{\PYZdq{}}\PY{p}{:} \PY{n}{rate}\PY{p}{,}
        \PY{l+s+s2}{\PYZdq{}}\PY{l+s+s2}{sample\PYZus{}size}\PY{l+s+s2}{\PYZdq{}}\PY{p}{:} \PY{n}{sample\PYZus{}size}\PY{p}{,}
        \PY{l+s+s2}{\PYZdq{}}\PY{l+s+s2}{moyenne\PYZus{}theorique}\PY{l+s+s2}{\PYZdq{}}\PY{p}{:} \PY{n}{theorical\PYZus{}average}\PY{p}{,}
        \PY{l+s+s2}{\PYZdq{}}\PY{l+s+s2}{moyenne\PYZus{}reel}\PY{l+s+s2}{\PYZdq{}}\PY{p}{:} \PY{n}{real\PYZus{}average}\PY{p}{,}
        \PY{l+s+s2}{\PYZdq{}}\PY{l+s+s2}{diff}\PY{l+s+s2}{\PYZdq{}}\PY{p}{:} \PY{n}{real\PYZus{}average} \PY{o}{\PYZhy{}} \PY{n}{theorical\PYZus{}average}\PY{p}{,}
        \PY{l+s+s2}{\PYZdq{}}\PY{l+s+s2}{peak}\PY{l+s+s2}{\PYZdq{}}\PY{p}{:} \PY{n}{np}\PY{o}{.}\PY{n}{max}\PY{p}{(}\PY{n}{active\PYZus{}players}\PY{p}{)}\PY{p}{,}
        \PY{l+s+s2}{\PYZdq{}}\PY{l+s+s2}{active\PYZus{}at\PYZus{}t}\PY{l+s+s2}{\PYZdq{}}\PY{p}{:} \PY{n}{active\PYZus{}at\PYZus{}t}\PY{p}{,}
        \PY{l+s+s2}{\PYZdq{}}\PY{l+s+s2}{arrivals}\PY{l+s+s2}{\PYZdq{}}\PY{p}{:} \PY{n}{arrivals}\PY{p}{,}
        \PY{l+s+s2}{\PYZdq{}}\PY{l+s+s2}{play\PYZus{}time}\PY{l+s+s2}{\PYZdq{}}\PY{p}{:} \PY{n}{play\PYZus{}time}
    \PY{p}{\PYZcb{}}\PY{p}{)}
\end{Verbatim}
\end{tcolorbox}

    \begin{tcolorbox}[breakable, size=fbox, boxrule=1pt, pad at break*=1mm,colback=cellbackground, colframe=cellborder]
\begin{Verbatim}[commandchars=\\\{\}]
\PY{n}{report} \PY{o}{=} \PY{l+s+s2}{\PYZdq{}}\PY{l+s+s2}{\PYZdq{}}
\PY{k}{for} \PY{n}{\PYZus{}r} \PY{o+ow}{in} \PY{n}{results}\PY{p}{:}
    \PY{n}{report} \PY{o}{+}\PY{o}{=} \PY{l+s+sa}{f}\PY{l+s+s2}{\PYZdq{}\PYZdq{}\PYZdq{}}\PY{l+s+s2}{\PYZsh{}\PYZsh{} }\PY{l+s+si}{\PYZob{}}\PY{n}{\PYZus{}r}\PY{p}{[}\PY{l+s+s2}{\PYZdq{}}\PY{l+s+s2}{rate}\PY{l+s+s2}{\PYZdq{}}\PY{p}{]}\PY{l+s+si}{\PYZcb{}}\PY{l+s+s2}{ branchements / minute}
\PY{l+s+s2}{    \PYZhy{} Temps de la simulation:}\PY{l+s+se}{\PYZbs{}t}\PY{l+s+si}{\PYZob{}}\PY{n}{\PYZus{}r}\PY{p}{[}\PY{l+s+s2}{\PYZdq{}}\PY{l+s+s2}{sample\PYZus{}size}\PY{l+s+s2}{\PYZdq{}}\PY{p}{]}\PY{l+s+si}{\PYZcb{}}\PY{l+s+s2}{ minutes}
\PY{l+s+s2}{    \PYZhy{} Moyennes de joueurs actifs:}
\PY{l+s+s2}{        \PYZhy{} Échantillon / réel: }\PY{l+s+si}{\PYZob{}}\PY{n}{\PYZus{}r}\PY{p}{[}\PY{l+s+s2}{\PYZdq{}}\PY{l+s+s2}{moyenne\PYZus{}reel}\PY{l+s+s2}{\PYZdq{}}\PY{p}{]}\PY{l+s+si}{:}\PY{l+s+s2}{.2f}\PY{l+s+si}{\PYZcb{}}
\PY{l+s+s2}{        \PYZhy{} Théorique:          }\PY{l+s+si}{\PYZob{}}\PY{n}{\PYZus{}r}\PY{p}{[}\PY{l+s+s2}{\PYZdq{}}\PY{l+s+s2}{moyenne\PYZus{}theorique}\PY{l+s+s2}{\PYZdq{}}\PY{p}{]}\PY{l+s+si}{:}\PY{l+s+s2}{.2f}\PY{l+s+si}{\PYZcb{}}
\PY{l+s+s2}{        \PYZhy{} Différence:         }\PY{l+s+si}{\PYZob{}}\PY{n}{\PYZus{}r}\PY{p}{[}\PY{l+s+s2}{\PYZdq{}}\PY{l+s+s2}{diff}\PY{l+s+s2}{\PYZdq{}}\PY{p}{]}\PY{l+s+si}{:}\PY{l+s+s2}{.2f}\PY{l+s+si}{\PYZcb{}}
\PY{l+s+s2}{    \PYZhy{} Peak de joueurs: }\PY{l+s+si}{\PYZob{}}\PY{n}{\PYZus{}r}\PY{p}{[}\PY{l+s+s2}{\PYZdq{}}\PY{l+s+s2}{peak}\PY{l+s+s2}{\PYZdq{}}\PY{p}{]}\PY{l+s+si}{\PYZcb{}}
\PY{l+s+s2}{    \PYZhy{} Joueurs actifs à différents instants:}
\PY{l+s+s2}{    }\PY{l+s+si}{\PYZob{}}\PY{l+s+s2}{\PYZdq{}}\PY{l+s+s2}{\PYZdq{}}\PY{o}{.}\PY{n}{join}\PY{p}{(}\PY{p}{[}\PY{l+s+sa}{f}\PY{l+s+s2}{\PYZdq{}}\PY{l+s+se}{\PYZbs{}t}\PY{l+s+s2}{\PYZhy{} t=}\PY{l+s+si}{\PYZob{}}\PY{n}{ti}\PY{l+s+si}{\PYZcb{}}\PY{l+s+s2}{ min : }\PY{l+s+si}{\PYZob{}}\PY{n}{n}\PY{l+s+si}{\PYZcb{}}\PY{l+s+se}{\PYZbs{}n}\PY{l+s+s2}{\PYZdq{}}\PY{+w}{ }\PY{k}{for}\PY{+w}{ }\PY{n}{ti}\PY{p}{,}\PY{+w}{ }\PY{n}{n}\PY{+w}{ }\PY{o+ow}{in}\PY{+w}{ }\PY{n}{\PYZus{}r}\PY{p}{[}\PY{l+s+s2}{\PYZdq{}}\PY{l+s+s2}{active\PYZus{}at\PYZus{}t}\PY{l+s+s2}{\PYZdq{}}\PY{p}{]}\PY{o}{.}\PY{n}{items}\PY{p}{(}\PY{p}{)}\PY{p}{]}\PY{p}{)}\PY{l+s+si}{\PYZcb{}}\PY{l+s+se}{\PYZbs{}n}\PY{l+s+s2}{\PYZdq{}\PYZdq{}\PYZdq{}}
\PY{n}{mo}\PY{o}{.}\PY{n}{md}\PY{p}{(}\PY{n}{report}\PY{p}{)}
\end{Verbatim}
\end{tcolorbox}



    \subsection{10 branchements / minute}\label{branchements-minute}

\begin{verbatim}
- Temps de la simulation:   10000 minutes
- Moyennes de joueurs actifs:
    - Échantillon / réel: 2810.91
    - Théorique:          2805.60
    - Différence:         5.31
- Peak de joueurs: 2946
- Joueurs actifs à différents instants:
    - t=10 min : 106
    - t=50 min : 517
    - t=500 min : 2840
\end{verbatim}

\subsection{67 branchements / minute}\label{branchements-minute-1}

\begin{verbatim}
- Temps de la simulation:   10000 minutes
- Moyennes de joueurs actifs:
    - Échantillon / réel: 18809.12
    - Théorique:          18797.52
    - Différence:         11.60
- Peak de joueurs: 19097
- Joueurs actifs à différents instants:
    - t=10 min : 662
    - t=50 min : 3358
    - t=500 min : 18863
\end{verbatim}

\subsection{100 branchements / minute}\label{branchements-minute-2}

\begin{verbatim}
- Temps de la simulation:   10000 minutes
- Moyennes de joueurs actifs:
    - Échantillon / réel: 28051.40
    - Théorique:          28056.00
    - Différence:         -4.60
- Peak de joueurs: 28393
- Joueurs actifs à différents instants:
    - t=10 min : 1000
    - t=50 min : 5033
    - t=500 min : 28192
\end{verbatim}



    \begin{tcolorbox}[breakable, size=fbox, boxrule=1pt, pad at break*=1mm,colback=cellbackground, colframe=cellborder]
\begin{Verbatim}[commandchars=\\\{\}]
\PY{k}{for} \PY{n}{\PYZus{}r} \PY{o+ow}{in} \PY{n}{results}\PY{p}{:}

    \PY{n}{P} \PY{o}{=} \PY{n}{generate\PYZus{}arrival\PYZus{}times}\PY{p}{(}\PY{n}{\PYZus{}r}\PY{p}{[}\PY{l+s+s1}{\PYZsq{}}\PY{l+s+s1}{rate}\PY{l+s+s1}{\PYZsq{}}\PY{p}{]}\PY{p}{,} \PY{n}{sample\PYZus{}size}\PY{p}{)}

    \PY{c+c1}{\PYZsh{} Histogram of generated inter\PYZhy{}arrival times}

    \PY{n}{plt}\PY{o}{.}\PY{n}{figure}\PY{p}{(}\PY{n}{figsize}\PY{o}{=}\PY{p}{(}\PY{l+m+mi}{10}\PY{p}{,}\PY{l+m+mi}{4}\PY{p}{)}\PY{p}{)}
    \PY{n}{plt}\PY{o}{.}\PY{n}{subplot}\PY{p}{(}\PY{l+m+mi}{1}\PY{p}{,}\PY{l+m+mi}{2}\PY{p}{,}\PY{l+m+mi}{1}\PY{p}{)}
    \PY{n}{plt}\PY{o}{.}\PY{n}{hist}\PY{p}{(}\PY{n}{P}\PY{p}{,} \PY{n}{bins}\PY{o}{=}\PY{l+m+mi}{50}\PY{p}{,} \PY{n}{density}\PY{o}{=}\PY{k+kc}{True}\PY{p}{,} \PY{n}{label}\PY{o}{=}\PY{l+s+s1}{\PYZsq{}}\PY{l+s+s1}{Generated P}\PY{l+s+s1}{\PYZsq{}}\PY{p}{)}
    \PY{c+c1}{\PYZsh{}plt.hist(\PYZus{}r[\PYZsq{}arrivals\PYZsq{}], bins=50, density=True)}
    \PY{n}{plt}\PY{o}{.}\PY{n}{title}\PY{p}{(}\PY{l+s+sa}{f}\PY{l+s+s2}{\PYZdq{}}\PY{l+s+s2}{Arrivées (λ=}\PY{l+s+si}{\PYZob{}}\PY{n}{\PYZus{}r}\PY{p}{[}\PY{l+s+s1}{\PYZsq{}}\PY{l+s+s1}{rate}\PY{l+s+s1}{\PYZsq{}}\PY{p}{]}\PY{l+s+si}{\PYZcb{}}\PY{l+s+s2}{/min)}\PY{l+s+s2}{\PYZdq{}}\PY{p}{)}

    \PY{n}{plt}\PY{o}{.}\PY{n}{subplot}\PY{p}{(}\PY{l+m+mi}{1}\PY{p}{,}\PY{l+m+mi}{2}\PY{p}{,}\PY{l+m+mi}{2}\PY{p}{)}
    \PY{n}{plt}\PY{o}{.}\PY{n}{hist}\PY{p}{(}\PY{n}{\PYZus{}r}\PY{p}{[}\PY{l+s+s1}{\PYZsq{}}\PY{l+s+s1}{play\PYZus{}time}\PY{l+s+s1}{\PYZsq{}}\PY{p}{]}\PY{p}{,} \PY{n}{bins}\PY{o}{=}\PY{l+m+mi}{50}\PY{p}{,} \PY{n}{density}\PY{o}{=}\PY{k+kc}{True}\PY{p}{)}
    \PY{n}{plt}\PY{o}{.}\PY{n}{title}\PY{p}{(}\PY{l+s+sa}{f}\PY{l+s+s2}{\PYZdq{}}\PY{l+s+s2}{Temps de jeu Q (μ=}\PY{l+s+si}{\PYZob{}}\PY{n}{MU\PYZus{}Q}\PY{l+s+si}{\PYZcb{}}\PY{l+s+s2}{ min)}\PY{l+s+s2}{\PYZdq{}}\PY{p}{)}

    \PY{n}{plt}\PY{o}{.}\PY{n}{tight\PYZus{}layout}\PY{p}{(}\PY{p}{)}
    \PY{n}{plt}\PY{o}{.}\PY{n}{show}\PY{p}{(}\PY{p}{)}
\end{Verbatim}
\end{tcolorbox}

    \begin{center}
    \adjustimage{max size={0.9\linewidth}{0.9\paperheight}}{figures/output_18_0.png}
    \end{center}
    { \hspace*{\fill} \\}

    \begin{center}
    \adjustimage{max size={0.9\linewidth}{0.9\paperheight}}{figures/output_18_1.png}
    \end{center}
    { \hspace*{\fill} \\}

    \begin{center}
    \adjustimage{max size={0.9\linewidth}{0.9\paperheight}}{figures/output_18_2.png}
    \end{center}
    { \hspace*{\fill} \\}

    \subsection{Conclusion de la
simulation}\label{conclusion-de-la-simulation}

C'est une belle simulation, mais malheureusement elle est très loin
d'une situation réel. En effet, elle assume un taux de branchement
constant à travers la journée, tandis que le taux change en fonction de
l'heure. Elle assume aussi une moyenne de temps de jeux constant, alors
qu'en réalité, les gens jouent plus les fins de semaines que les
journées de travail.

Tout ceci ne prend pas en compte le plus gros problème potentiel.
Lorsqu'un jeux est mis sur le marché, on observe une fonction
décroissante du nombre de joueur à travers le temps, avec un énorme
montant de joueurs au tout début de la mise en marché. S'ils utilisent
des moyennes pour quantifié leurs besoins en serveurs, la réalité vas
vite les rattrapés lors du lancement du jeux.

Ils ont besoin au minimum du double de l'infrastructure prévus pour les
premières semaines de la mises en marché, avec possibilités
d'aggrendissements temporaires lors d'évènements saisoniers ou des
periodes de fêtes nationaux.

En bref, c'est une bonne simulation de moyenne, mais ça ne devrait pas
être la seule métrique de prise de décisions pour l'infrastructure de
serveurs.
